\documentclass[conference]{IEEEtran}
\usepackage{cite}
\usepackage{amsmath,amssymb,amsfonts}
\usepackage{algorithmic}
\usepackage{graphicx}
\usepackage{textcomp}
\usepackage{xcolor}
\usepackage{url}
\usepackage{hyperref}
\usepackage{float}
\usepackage{listings}
\lstset{literate=%
    {├}{{\textSFviii}}1
    {─}{{\textSFx}}1
    {│}{{\textSFxi}}1
    {└}{{\textSFii}}1
}
\lstdefinelanguage{json}{
    basicstyle=\ttfamily\small,
    numbers=left,
    numberstyle=\scriptsize,
    stepnumber=1,
    numbersep=8pt,
    showstringspaces=false,
    breaklines=true,
    frame=lines,
    string=[s]{"}{"},
    stringstyle=\color{blue},
    comment=[l]{:},
    commentstyle=\color{black},
}

\lstdefinelanguage{TypeScript}{
    keywords={abstract, any, as, break, case, catch, class, const, continue, debugger, declare, default, delete,
    do, else, enum, export, extends, false, finally, for, from, function, get, if, implements, import, in, 
    instanceof, interface, let, module, namespace, new, null, number, of, package, private, protected, public,
    readonly, require, return, set, static, string, super, switch, symbol, this, throw, true, try, type, typeof,
    var, void, while, with, yield},
    keywordstyle=\color{blue}\bfseries,
    ndkeywords={boolean, number, string, any, void, never},
    ndkeywordstyle=\color{purple}\bfseries,
    identifierstyle=\color{black},
    sensitive=false,
    comment=[l]{//},
    morecomment=[s]{/*}{*/},
    commentstyle=\color{gray}\ttfamily,
    stringstyle=\color{red}\ttfamily,
    morestring=[b]',
    morestring=[b]",
    basicstyle=\ttfamily\small,
    numbers=left,
    numberstyle=\scriptsize,
    stepnumber=1,
    numbersep=8pt,
    showstringspaces=false,
    breaklines=true,
    frame=lines,
}

\setlength{\textfloatsep}{3pt plus 1.0pt minus 1.0pt} 
\setlength{\floatsep}{3pt plus 1.0pt minus 1.0pt} 
\setlength{\intextsep}{3pt plus 1.0pt minus 1.0pt}
\setlength{\abovecaptionskip}{3pt}
\setlength{\belowcaptionskip}{3pt}

\renewcommand{\figurename}{Fig.}
\renewcommand{\tablename}{Tabel}
\renewcommand{\refname}{Bibliografie}

\begin{document}

\title{Proiectarea și implementarea unei platforme web pentru managementul fluxurilor studențești utilizând Angular și PHP}

\author{
    \IEEEauthorblockN{Student: Gherguț Mădălin}
    \IEEEauthorblockA{Grupa: 1405A \\
        \textit{Facultatea de Automatică și Calculatoare} \\
        \textit{Universitatea Tehnică ``Gheorghe Asachi'' din Iași}}
    \textit{email: madalin.ghergut@student.tuiasi.ro}
    \and
    \IEEEauthorblockN{Îndrumător: Ș.l.dr.ing. Cătălin Brăescu}
    \IEEEauthorblockA{\textit{Facultatea de Automatică și Calculatoare} \\
        \textit{Universitatea Tehnică ``Gheorghe Asachi'' din Iași}}
    \textit{email: florin-catalin.braescu@academic.tuiasi.ro}

}

\maketitle

\begin{abstract}
    Raportul documentează studiul si construirea unei platforme web pentru managementul studenților și urmărirea prezenței pentru semestrul I.
    Activitățile principale au cuprins: (1) însuşirea conceptelor fundamentale privind arhitectura SPA și API-urile RESTful;
    (2) familiarizarea cu Angular 19 --- componente, servicii, rutare; (3) documentarea creării API-uri REST în PHP cu gestionarea
    securizată a cererilor HTTP și conexiuni MySQL; (4) proiectarea schemei bazei de date relaționale cu diagrama E-R;
    (5) design preliminar al interfeței utilizator. Tehnologiile principale sunt Angular 19.2.0 pentru frontend,
    PHP cu PDO pentru backend, și MySQL pentru persistență. Documentul prezintă conceptele teoretice fundamentale
    (Client-Server, HTTP/REST, JSON, Reactive Programming, DI), arhitectura aplicației, patternuri de design și
    planning pentru Semestrul II inclusiv testare comprehensivă și funcționalități avansate.
\end{abstract}

\begin{IEEEkeywords}
    Angular 19, REST API, PHP, Student Management System, Attendance Tracking, Single Page Application, MySQL, Web Application,
    Database Design, Role-Based Access Control.
\end{IEEEkeywords}

\section{Introducere}
În era digitalizării instituțiilor de învățământ, transformarea proceselor administrative și academice prin intermediul tehnologiei
informației a devenit o necesitate. Sistemele tradiționale de urmărire a prezenței, bazate pe registre fizice și calcule manuale,
se confruntă cu limitări evidente: proastă scalabilitate, erori umane frecvente, dificultăți în generarea rapoartelor.
Dezvoltarea aplicațiilor web moderne a deschis noi oportunități pentru modernizarea acestor procese, permitând instituțiilor să ofere soluții eficiente,
transparente și accesibile din orice loc și oricând.

În contextul acestui proiect, am identificat necesitatea creării unei platforme digitale care să consolideze toate aspectele management-ului studenților
și urmărirea prezenței într-o singură aplicație cohesivă. Tehnologiile web moderne, în special paradigma Single Page Application (SPA)
combinată cu API-uri RESTful, oferă arhitectura ideală pentru a realiza această viziune. Alegerea Angular 19 ca
framework frontend și PHP ca limbaj backend reflectă un echilibru între maturitate tehnologică, disponibilitate de resurse și
ușurință de întreținere pe termen lung. Conform sondajului Stack Overflow Developer Survey 2025, Angular este folosit de 18.2\%
dintre dezvoltatorii web profesioniști~\cite{StackOverflow2025}, demonstrând adoptarea largă și comunitatea activă în jurul acestui framework.

Obiectivul principal al proiectului este crearea unei platforme care să
faciliteze:

\begin{itemize}
    \item Gestionarea centralizată a studenților, profesorilor și grupelor de curs
    \item Urmărirea automată a prezenței studenților la sesiuni
    \item Managementul laboratoarelor și al programului de curs
    \item Control de acces bazat pe roluri (admin, profesor, student)
    \item Vizualizare în timp real a datelor și rapoarte
    \item Scalabilitate pentru instituții cu sute de studenți și profesori
\end{itemize}

Arhitectura aleasă este de tip client-server, cu o interfață frontend modulară
construită în Angular 19~\cite{Angular2024} și un backend RESTful implementat
în PHP cu bază de date MySQL~\cite{MySQL2024}.

\section{Arhitectura Sistemului și Tehnologii}
O înțelegere solidă a arhitecturii sistemului și a conceptelor fundamentale
este esențială pentru proiectarea unei aplicații web moderne și scalabile. În
această secțiune, prezentăm conceptele teoretice pe care se construiește
întregul proiect.

\subsection{Arhitectura Client-Server}
Modelul client-server este o paradigmă arhitectural fundamental în tehnologie
web. În acest model:

\begin{itemize}
    \item Client: O aplicație (browser, mobil, desktop) care inițiază cererile către
          server
    \item Server: O aplicație care primește cererile, le procesează și trimite răspunsuri
    \item Comunicație: Se realizează prin protocoale standardizate (HTTP/HTTPS)
\end{itemize}

Pentru proiectul nostru:
\begin{itemize}
    \item Client: Aplicația Angular care rulează în browserul utilizatorului
    \item Server: Backend-ul PHP care gestiunează baza de date și logica de business
    \item Protocol: HTTP/REST pentru comunicare
\end{itemize}

Avantajele acestui model includ:
\begin{itemize}
    \item Separarea responsabilităților: Frontend și backend pot evolua independent
    \item Scalabilitate: Serverul poate gestiona multiple clienți
\end{itemize}

Comunicarea client-server se realizează conform diagramei prezentate în Figura
~\ref{fig:client_server_arch}, unde clientul trimite cereri HTTP iar serverul
răspunde cu date în format JSON.

\begin{figure}[H]
    \centering
    \includegraphics[width=0.48\textwidth]{images/arhitecture.png}
    \caption{Arhitectura de ansamblu a sistemului, ilustrând separarea între clientul Single Page Application (Angular)
        și serverul backend (PHP și MySQL). Comunicarea se realizează asincron prin cereri HTTP, utilizând formatul JSON pentru transferul datelor.}%
    \label{fig:client_server_arch}
\end{figure}

\subsection{Protocolul HTTP și Metodele RESTful}
HTTP (Hypertext Transfer Protocol) este protocolul fundamental al web-ului.
Este un protocol stateless, ceea ce înseamnă că fiecare cerere este
independentă și nu ține minte starea cererilor anterioare.

Metodele HTTP principale utilizate în API-urile RESTful sunt:

\begin{itemize}
    \item GET: Recuperează resurse fără a modifica serverul (idempotent)
    \item POST: Creează o nouă resursă pe server
    \item PUT: Actualizează complet o resursă existentă
    \item PATCH: Actualizare parțială a unei resurse
    \item DELETE: Șterge o resursă
\end{itemize}

Exemplu de cereri HTTP în contextul proiectului:
\begin{itemize}
    \item \texttt{GET /api/courses} --- Listă toate cursurile
    \item \texttt{POST /api/students} --- Creează un student nou
    \item \texttt{PUT /api/attendance/1} --- Actualizează un record de prezență
    \item \texttt{DELETE /api/courses/5} --- Șterge un curs
\end{itemize}

\subsection{JSON --- Formatul de Schimb de Date}

JSON (JavaScript Object Notation) este un format ușor și text-based pentru
reprezentarea datelor structurate. Este limbajul standard pentru comunicarea în
API-urile RESTful moderne.

Exemplu de răspuns JSON din aplicația noastră:

\begin{lstlisting}[language=json]
{
  "id": 1,
  "name": "Structuri de Date",
  "teacher_id": 5,
  "description": "Introducere in structuri...",
  "created_at": "2024-01-15T10:30:00Z"
}
\end{lstlisting}

Avantajele JSON:
\begin{itemize}
    \item Ușor de citit pentru oameni și ușor de parsat de mașini
    \item Suport nativ în JavaScript și TypeScript
    \item Format text, independent de platformă
    \item Suport pentru tipuri de date complexe (obiecte, matrici)
\end{itemize}

\subsection{Modelul de Date Relațional}

Baza de date a proiectului utilizează modelul relațional, propus de E.F. Codd
în 1970~\cite{Codd1970}. În acest model:

\begin{itemize}
    \item Datele sunt organizate în tabele (relații)
    \item Fiecare tabel are coloane (atribute) și rânduri (înregistrări)
    \item Chei primare identifică unic fiecare înregistrare
    \item Chei străine stabilesc relații între tabele
    \item Constrângeri de integritate asigură consistența datelor
\end{itemize}

Figura~\ref{fig:rdb_model} ilustrează conceptele fundamentale ale modelului
relațional, prezentând modul în care datele sunt structurate în tabele cu chei
primare și străine, precum și relațiile dintre entități.

\begin{figure}[H]
    \centering
    \includegraphics[width=0.48\textwidth]{images/RDB.png}
    \caption{Modelul de date relațional. Ilustrează conceptele fundamentale: tabele,
        atribute, chei primare (PK), chei străine (FK) și relațiile dintre entități.}%
    \label{fig:rdb_model}
\end{figure}

\subsection{Programarea Reactivă și Observables}

Programarea reactivă este un paradim de programare care se concentrează pe
fluxurile de date asincrone și propagarea schimbărilor. În Angular, aceasta se
realizează prin RxJS Observables~\cite{RxJS2024}.

Un Observable este o reprezentare a unei serii de valori care pot fi emise în
timp, permițându-ți să te abonezi pentru a primi aceste valori pe măsură ce
sunt disponibile.

\begin{itemize}
    \item Asincron: Valorile sunt emise în viitor, nu imediat
    \item Lazy: Observable-ul nu se execută până când nu este abonat
    \item Cancellable: Abonamentul poate fi anulat în orice moment
\end{itemize}

Exemplu de caz de folosință: Obținerea listei de studenți de la server și
afișandule în consolă:

\begin{lstlisting}[language=TypeScript]
this.studentService.getStudents()
  .subscribe(students => {
    console.log(students);
  });
\end{lstlisting}

\subsection{Injectarea dependențelor}

Injectarea dependențelor este un model de proiectare care promovează decuplarea
componentelor. În loc ca o clasă să creeze obiectele de care are nevoie,
aceasta le primește de la exterior.

Beneficii ale injectării dependențelor:
\begin{itemize}
    \item Reutilizabilitate: Componentele devin mai independente și reutilizabile
    \item Modularitate: Codul este mai bine organizat
    \item Flexibilitate: Modificări ușoare în implementare fără afectarea consumatorilor
\end{itemize}

Angular are un sistem de injectare a dependențelor încorporat care gestionează
automatica injectarea dependențelor în componente și servicii.

\section{Arhitectura Aplicațiilor SPA și API-uri RESTful}

\subsection{Concepte SPA și Avantaje în Contextul Actual}

O aplicație Single Page Application (SPA) este o aplicație web care se încarcă
o singură dată în browser și apoi interacționează cu utilizatorul prin
intermediul JavaScript, fără reîncărcări complete de pagină~\cite{SPABook2024}.
În contextul actual, SPA-urile au devenit standard în dezvoltarea web modernă
din cauza avantajelor semnificative pe care le oferă.

Caracteristicile principale ale unei SPA includ:

\begin{itemize}
    \item Interacțiune fluidă: Interfața răspunde instantaneu, oferind o experiență
          similară aplicațiilor clasice instalate pe calculator.
    \item Actualizare dinamică: Datele se transferă în fundal, permițând reîmprospătarea
          informațiilor fără reîncărcarea paginii.
    \item Navigare locală: Trecerea între secțiuni se face direct în browser, reducând
          solicitarea serverului și timpul de așteptare.
    \item Gestiune locală a datelor: Stocarea informațiilor în memoria browserului
          permite accesul rapid și stabilitatea în cazul unei conexiuni slabe.
    \item Arhitectură modulară: Separarea clară între interfață și logica de procesare a
          datelor facilitează întreținerea sistemului.
\end{itemize}

\subsection{API-uri RESTful}

Arhitectura REST (Representational State Transfer) definește un set de
principii pentru proiectarea serviciilor web~\cite{AWSRESTful2024}. În
contextul proiectului, backend-ul expone endpoint-uri care respectă standardele
REST:

\begin{itemize}
    \item Resurse: Entități principale (utilizatori, cursuri, prezențe)
    \item Operații HTTP: GET (citire), POST (creare), PUT (actualizare), DELETE
          (ștergere)
    \item Reprezentare JSON: Schimb de date structurat
    \item Coduri de stare HTTP: Comunicarea rezultatelor operațiilor prin coduri
          standardizate:
          \begin{itemize}
              \item 200 OK: Cererea a fost procesată cu succes
              \item 201 Created: Resursa a fost creată cu succes
              \item 400 Bad Request: Cererea este invalidă sau malformată
              \item 404 Not Found: Resursa solicitată nu a fost găsită
              \item 500 Internal Server Error: Eroare pe server la procesarea cererii
          \end{itemize}
\end{itemize}

\section{Proiectare și Implementare}

\subsection{Framework Angular 19 --- Implementare}

Angular este un framework TypeScript open-source, dezvoltat și întreținut de
Google, destinat construirii aplicațiilor web scalabile și
performante~\cite{Angular2024}. Angular oferă o soluție completă pentru
dezvoltarea SPA-urilor, incluzând:

\begin{itemize}
    \item Sistem de componente: Arhitectură bazată pe componente reutilizabile
    \item Rutare: Navigator integrat pentru aplicații multi-pagină (pe client)
    \item Gestionarea formularelor: Suport pentru validare complexă și legare
          bidirecțională a datelor
    \item Servicii: Injecție de dependență pentru cod modular și testabil
    \item HTTP Client: Comunicație ușoară cu backend-ul
    \item RxJS Reactive Extensions: Programare reactivă și gestionarea evenimentelor
\end{itemize}

\subsection{Componente și Servicii}

Componentele Angular sunt blocuri reutilizabile ce încapsulează logica,
template-ul și stilurile. Serviciile furnizează funcționalități partajate și
comunicația cu API-ul. Relația dintre componente și servicii este ilustrată în
Figura~\ref{fig:angular_internal_arch}.

\begin{figure}[H]
    \centering
    \includegraphics[width=0.48\textwidth]{images/angular-architecture.png}
    \caption{Arhitectura modulară a frontend-ului Angular. Sunt ilustrate ierarhia componentelor (gestionată de Router),
        separarea logicii de prezentare de cea de business prin intermediul serviciilor, și integrarea acestora folosind
        mecanismul de injectare a dependențelor (Dependency Injection).}%
    \label{fig:angular_internal_arch}
\end{figure}

\subsection{Rutare și Managementul Stării}

Aplicația utilizează sistemul de rutare Angular pentru navigarea între
secțiuni, configurat în fișierul \texttt{app.routes.ts}. Rutele principale ale
aplicației includ:

\begin{itemize}
    \item \texttt{/login} --- Autentificare și validare credențiale
    \item \texttt{/dashboard} --- Panou principal cu statistici generale
    \item \texttt{/courses} --- Gestionare cursuri (vizualizare, creare, editare)
    \item \texttt{/students} --- Administrare studenți (operații CRUD)
    \item \texttt{/attendance} --- Înregistrare și monitorizare prezență
\end{itemize}

\subsection{Arhitectura Backend}

Backend-ul a fost construit în PHP~\cite{PHP2024} cu utilizarea PDO (PHP Data
Objects) pentru conexiunile securizate la baza de date. Structura proiectului
organizează endpoint-urile după resurse, cu API-uri CRUD pentru gestionarea
studenților, cursurilor, grupelor și prezențelor:

\begin{lstlisting}[language=bash]
api/
|-- auth/
|   |-- login.php
|   +-- verify.php
|-- courses/
|   |-- getAll.php
|   |-- create.php
|   |-- update.php
|   +-- delete.php
|-- students/
|   |-- getAll.php
|   |-- create.php
|   |-- update.php
|   +-- delete.php
+-- attendance/
    |-- save.php
    |-- getByStudent.php
    +-- getStatsByCourse.php
\end{lstlisting}

\subsection{Schema Bazei de Date}

Baza de date a fost proiectată urmând modelul relațional, cu entități
principale și relații dintre ele:

\begin{itemize}
    \item Users: Credențiale de autentificare pentru toți utilizatorii
    \item Teachers: Informații profesori (ref. Users)
    \item Students: Informații studenți și grup (ref. Users, Groups)
    \item Groups: Grupe de studenți
    \item Courses: Cursuri universitare (ref. Teachers)
    \item GroupCourses: Relație M:N între grupe și cursuri
    \item Sessions: Sesiuni de curs (ref. GroupCourses)
    \item Laboratories: Laboratoare per curs (1--14 per curs)
    \item Attendance: Prezența studenților la sesiuni (ref. Students, Sessions)
\end{itemize}

\begin{figure}[H]
    \centering
    \includegraphics[width=0.48\textwidth]{images/database-er-diagram.png}
    \caption{Diagrama Entitate-Relație (ERD) a bazei de date. Schema definește entitățile și relațiile de
        cardinalitate necesare pentru gestionarea studenților, cursurilor, sesiunilor și prezenței.}%
    \label{fig:er_diagram}
\end{figure}

\subsection{Relații între Tabele}

Schema este normalizată la a treia formă normală (3NF), cu constrângeri de
integritate referențială care asigură consistența datelor:

\begin{itemize}
    \item Un profesor are mai multe cursuri (1:N)
    \item Un curs are mai multe sesiuni (1:N)
    \item O grupă participă la mai multe cursuri prin GroupCourses (M:N)
    \item Un student aparține unei grupe (N:1)
    \item O sesiune are mai multe prezențe (1:N)
    \item Tabelul GroupCourses permite unui curs să fie oferit mai multor grupe
\end{itemize}

\subsection{Interfața Utilizator și Funcționalități}

Aplicația include interfețe dedicate pentru fiecare rol. Secțiunile de monitorizare 
prezențe este destinată studenților, iar interfața 
de creare sesiuni laborator este destinată profesorilor.


\subsubsection{Autentificare și Securitate}

Sistemul implementează autentificare securizată care verifică credențialele
utilizatorului și atribuie roluri corespunzătoare (student, profesor, admin).
Utilizatorii se autentifică cu email și parolă, iar sistemul redirecționează
automat către dashboard-ul personalizat în funcție de rolul lor.
\begin{figure}[H]
    \centering
    \includegraphics[width=0.48\textwidth]{images/login-page.png}
    \caption{Pagina de autentificare a platformei. Interfață simplă și intuitivă cu validare pe client.}%
    \label{fig:login_page_alt}
\end{figure}

\subsubsection{Monitorizare Prezențe --- Situație Generală}

În partea superioară a dashboard-ului, sistemul oferă studenților o vizualizare agregată a situației
lor de prezență la toate cursurile. Această secțiune afișează statistici generale prin grafice pie
chart pentru fiecare materie, procentaje de prezență și numărul total de sesiuni la care au participat.

\begin{figure}[H]
    \centering
    \includegraphics[width=0.48\textwidth]{images/attendace-overview.png}
    \caption{Secțiunea situație generală prezențe din dashboard-ul studentului. Afișează grafice pie chart pentru fiecare curs, prezentând vizual situația prezenței (prezent/absent).}%
    \label{fig:attendance_overview}
\end{figure}

\subsubsection{Monitorizare Prezențe --- Lista Cronologică}

În partea inferioară a dashboard-ului, studenții pot consulta lista detaliată
și cronologică a tuturor sesiunilor de laborator la care au participat sau
au fost absenți.

\begin{figure}[H]
    \centering
    \includegraphics[width=0.48\textwidth]{images/cronological-overview.png}
    \caption{Lista cronologică a prezențelor din dashboard-ul studentului.}%
    \label{fig:attendance_list}
\end{figure}

\subsubsection{Interfața Profesor --- Creare Sesiune Laborator}

Profesorii dispun de o interfață pentru crearea și gestionarea sesiunilor de 
laborator. Funcționalitățile principale includ selectarea cursului, grupei și 
laboratoarelor specifice, precum și 
flexibilitatea de a adăuga studenți din alte grupe pentru recuperări sau situații speciale.

\begin{figure}[H]
    \centering
    \includegraphics[width=0.48\textwidth]{images/creating-laboratory.png}
    \caption{Interfața de creare sesiune laborator pentru profesori.}%
    \label{fig:teacher_create_session}
\end{figure}

\subsubsection{Interfața Profesor --- Lista Studenți}

Profesorii au acces la o interfață de vizualizare a listei complete de studenți, 
cu posibilitatea de filtrare după grupă și paginare pentru navigare eficientă.

\begin{figure}[H]
    \centering
    \includegraphics[width=0.48\textwidth]{images/student-list-profesor-view.png}
    \caption{Interfața de vizualizare studenți pentru profesori cu filtrare după grupă și paginare.}%
    \label{fig:student-list-profesor-view}
\end{figure}

\section{Direcții de Dezvoltare}

\subsection{Planul pentru Semestrul II}

Extinderi și îmbunătățiri planificate pentru Semestrul II:

\begin{itemize}
    \item Modelare baza de date pentru a respecta cursurile opționale
    \item Aprofundare CLI Angular
    \item Teste unitare cu Jasmine și Karma
    \item Teste end-to-end cu Cypress
    \item Control de acces bazat pe roluri (RBAC)
    \item Route guards și protejarea rutelor
    \item Optimizare frontend: lazy loading, change detection, caching
    \item Optimizare backend: database indexing, query optimization
    \item Documentare deployment pe un server local
    \item Imbunatarire interfata utilizator
\end{itemize}
\begin{thebibliography}{99}

    \bibitem{Angular2024}
    Google, \textit{Angular v19 Overview}, https://v19.angular.dev/overview, 2024.

    \bibitem{RxJS2024}
    RxJS Team, \textit{RxJS Documentation}, https://rxjs.dev/, 2024.

    \bibitem{PHP2024}
    PHP Group, \textit{PHP Manual}, https://www.php.net/manual/, 2024.

    \bibitem{MySQL2024}
    Oracle, \textit{MySQL Documentation}, https://dev.mysql.com/doc/, 2024.

    \bibitem{Codd1970}
    E.F. Codd, \textit{A Relational Model of Data for Large Shared Data Banks},
    Communications of the ACM, vol. 13, no. 6, pp. 377--387, 1970.
    doi:10.1145/362384.362685.

    \bibitem{AWSRESTful2024}
    Amazon Web Services, \textit{What is a RESTful API?},
    https://aws.amazon.com/what-is/restful-api/, 2024.

    \bibitem{SPABook2024}
    Miško Hevery and others, \textit{Single Page Application Book},
    https://singlepageappbook.com/, 2024.

    \bibitem{StackOverflow2025}
    Stack Overflow, \textit{Stack Overflow Developer Survey 2025}, https://survey.stackoverflow.co/2025/, 2025.

\end{thebibliography}

\end{document}
